% ================================================================
%  conteudo/justificativa.tex
% ================================================================

\section{Justificativa}

A aviação civil responde por cerca de 915 milhões de toneladas de
CO\textsubscript{2} emitidas anualmente, e projeções indicam que a demanda
por transporte aéreo dobrará até 2050 (IATA, 2023). Diferentemente dos
setores rodoviário e marítimo, onde a eletrificação avança rapidamente, as
aeronaves de longo curso são fisicamente limitadas pela densidade energética
dos combustíveis líquidos, tornando o SAF a principal alternativa de
descarbonização disponível no horizonte tecnológico atual.

O Brasil ocupa posição estratégica nesse cenário: o país possui extensa
faixa costeira e regiões de alta irradiância solar, condições favoráveis
ao cultivo de microalgas em larga escala, além de parque industrial
consolidado com geração de efluentes que podem substituir insumos de
fertilizantes nitrogenados e fosfatados no meio de cultivo microalgal
(MATA et al., 2010).

A espécie \textit{Chlorella fusca}, embora menos estudada que \textit{C.
vulgaris} e \textit{Nannochloropsis} sp., apresenta tolerância a variações
de pH, salinidade e carga orgânica, tornando-a candidata promissora para
o uso de efluentes industriais como substrato. A linha de pesquisa do
presente trabalho insere-se no eixo de biorrefinaria de microalgas --
que combina produção de energia, tratamento de resíduos e geração de
co-produtos de alto valor agregado --, tema de crescente interesse na
comunidade científica e no setor produtivo (MATA et al., 2010;
CHISTI, 2007).

Do ponto de vista de novidade científica, a literatura disponível sobre
\textit{C. fusca} é escassa em relação ao perfil lipídico obtido em
diferentes condições de cultivo e ao rendimento HEFA calculado a partir
de dados experimentais reais da espécie. A integração de simulação de
processo (cálculo estequiométrico HEFA) com dados experimentais obtidos
por GC-MS e a análise orçamentária com substituição de efluente industrial
constituem a contribuição original deste trabalho ao estado da arte.

A equipe possui histórico em análises de biomassa microbiana, tendo
desenvolvido metodologia validada de curva padrão para quantificação de
biomassa de \textit{Saccharomyces cerevisiae} por turbidimetria e
gravimetria, fornecendo base experimental consolidada para as etapas de
monitoramento cinético previstas neste projeto.

